\section{Ask the People It Happens To}
\label{sec:8.2}
Democratizing AI development means building the people affected by a system into how it gets built, not just consulting them once it already exists. A design team, a regulator, a user, and a community already carrying a system's costs each see a different part of what it does, and leaving any of them out means the decision gets made by whoever is left in the room.

The mechanisms for doing that are well established: cross-sector funding, inclusive design workshops held while a system is still a prototype, ethics advisory boards, and development conducted with a community instead of for it. Assembling that list is easy, and the assembling hides what separates the items on it. Three of the four make an AI system better informed. One of them can stop it. An advisory board is the only mechanism here that carries anything like a veto, and it carries one to the degree that the board can survive using it. A workshop that surfaces a serious objection has done its work and has no further move available. A board that surfaces the same objection and is overruled has demonstrated something about the institution that convened it. A board's worth is therefore measured by what happens the first time it says no.

What governance adds to development is machinery for bringing public input to bear on a decision at a single point in time instead of through an ongoing relationship.

\begin{enumerate}
  \item Citizen assemblies. A panel of citizens selected by random, stratified sampling deliberates on a question over an extended period and produces recommendations grounded in that deliberation rather than in a snap poll. Ireland's Citizens' Assembly has used this format to shape national recommendations on abortion law and on climate policy — questions carrying the same mix of technical complexity and moral stakes that AI governance has \autocite{citizensassembly2016ireland}. The format has since been pointed at AI directly. Forty randomly selected Americans, stratified across age, race, education, employment, party and region, spent eight days on high-risk AI in October 2023 \autocite{cndp2023assembly}. Taiwan's digital ministry convened a citizens' deliberative assembly in March 2024 on using AI to support information integrity, its participants drawn from across the country by stratified random sampling \autocite{moda2024integrity}. And forty randomly selected residents of French-speaking Switzerland spent two weekends in November 2025 producing twenty proposals on AI governance \autocite{epfl2025assembly}. The Ada Lovelace Institute's Citizens' Biometrics Council shows what the format needs from its question: fifty members of the UK public were asked not what they thought about AI, a question almost nobody can answer usefully, but about facial recognition — a thing with a location, an operator, and a consequence \autocite{adalovelaceinstitute2021biometrics}.
  \item Online platforms for public input. A standing online consultation reaches people a physical assembly cannot, at the cost of the depth that extended facilitated deliberation provides. The European Commission's consultation on its AI White Paper drew more than 1,200 responses from citizens, companies, and civil-society groups, and fed into the legislative proposal that became the EU AI Act \autocite{europeancommission2020whitepaper}.
  \item AI ethics committees. An ethics board answerable only to the company that created it is only as independent as that company is willing to let it be. DeepMind Health set up an outside panel in 2016 to review its work with the UK's National Health Service and publish annual reports on its findings; Google disbanded the panel in 2019, after members had raised concerns about how much access to information, and how much independence from the company, they actually had \autocite{fisher2019deepmind}. The panel's short life is the governance lesson, and it is the answer to what happens the first time a board says no: the board said something the company did not want to hear, and the company's available response was to stop having a board. Nothing in the arrangement had made that response costly. A review body is worth what it can enforce, and this one could enforce publication of its own concerns right up until the moment it could not. That is section~\ref{sec:3.5}'s halt in institutional form: what can be switched off cannot hold a line.
\end{enumerate}

\runin{What none of these mechanisms does}
There is a limit these mechanisms share. Every governance mechanism proposed in this chapter — citizen assemblies, standing consultations, inclusive design workshops, participatory design, the vTaiwan model — sits on the same side of a line. All of them make public preference more legible, more informed, and more actionable. Not one of them protects anybody from it. They are aggregation instruments, and a better aggregation instrument produces a more faithful rendering of what the public wants, including when what the public wants is something done to a minority of it.

This observation is not an objection to any of the mechanisms in particular. Preference aggregation is a real requirement and these are good ways to meet it. Meeting this requirement \emph{completely} would leave a central problem untouched. A perfectly deliberative, perfectly informed, perfectly representative process that concludes an entire demographic should be surveilled, deported, imprisoned, or otherwise scapegoated has produced a faultless output by every criterion in this chapter.

The argument about learning from human feedback gets close to this in noticing that treating all disagreement as equally worth representing goes wrong when one side's position is the product of coercion. But that version of the trouble still locates it in the input. The structural version of the claim is stronger and does not depend on anyone's preference being contaminated at all: some protections have to hold against the aggregate regardless of how the aggregate came out, and no amount of improvement to the aggregating machinery can supply them.

What supplies those protections is a different kind of machinery: enforceable rights that do not depend on being popular, a court willing to rule against a government that holds a legitimate majority, limits that whoever won the last election cannot repeal, and standing for people who are affected without being enfranchised. That half matters especially where a democracy is behaving most democratically.
